%---------------------------------------------------------------------
\section{Supplementation}
\label{sec:supplements}

%---------------------------------------------------------------------
\subsection*{Overview}

No supplement covered in this section has been tested in a climbing-specific RCT with a functional primary outcome (redpoint grade, route completion, pulley injury incidence, or return-to-climbing time). All evidence tiers for climbing applications are therefore extrapolated from adjacent populations or from surrogate-biomarker protocols. Claims confirmed across multiple independent literature queries are reported without qualification; single-source claims carry an explicit footnote and a \toverify flag where applicable.

%---------------------------------------------------------------------
\subsection{Study Summary Table}
\label{subsec:supptable}

\begin{longtable}{>{\raggedright\arraybackslash}p{2.8cm}
                  >{\raggedright\arraybackslash}p{2.4cm}
                  >{\raggedright\arraybackslash}p{2.6cm}
                  >{\raggedright\arraybackslash}p{2.6cm}
                  >{\raggedright\arraybackslash}p{2.0cm}
                  >{\raggedright\arraybackslash}p{1.6cm}
                  >{\raggedright\arraybackslash}p{1.0cm}
                  >{\raggedright\arraybackslash}p{1.2cm}}
\caption{Peer-reviewed intervention studies cited in this section.
         P1NP = amino-terminal propeptide of type I procollagen (collagen synthesis surrogate).
         VISA-A = Victorian Institute of Sport Assessment -- Achilles tendinopathy score.
         OSS = Oxford Shoulder Score; SPADI = Shoulder Pain and Disability Index.
         g = Hedges $g$; $d$ = Cohen's $d$.
         Endpoint: S = surrogate; F = functional.
         ES not reported = authors did not publish standardised effect size or 95\,\% CI.}\label{tab:suppstudy}\\
\toprule
\textbf{Study} & \textbf{Design / n} & \textbf{Supplement / dose} & \textbf{Primary outcome} & \textbf{Effect size} & \textbf{Endpoint} & \textbf{Qual.} & \textbf{Industry} \\
\midrule
\endfirsthead
\toprule
\textbf{Study} & \textbf{Design / n} & \textbf{Supplement / dose} & \textbf{Primary outcome} & \textbf{Effect size} & \textbf{Endpoint} & \textbf{Qual.} & \textbf{Industry} \\
\midrule
\endhead
\midrule
\multicolumn{8}{r}{\small\itshape Continued on next page}\\
\endfoot
\bottomrule
\endlastfoot

Shaw et al.\ (2017).
\textit{Am J Clin Nutr} 105(1):136--143. \cite{SLBR+17}
&
Randomised double-blind crossover; n=8 healthy males, age $21{\pm}0.5$\,yr; no tendon pathology. Washout 7\,days.
&
5\,g or 15\,g gelatin $+$ 48\,mg vitamin~C vs.\ placebo; single dose 1\,h before 6\,min rope-skipping.
&
Serum P1NP (AUC, 4\,h); mechanical properties of in-vitro engineered ligament constructs incubated with participant serum.
&
15\,g: ${\approx}2{\times}$ P1NP AUC vs.\ placebo ($p{<}0.05$); in-vitro ligament stiffness and collagen content $\uparrow$ significantly. Formal Cohen's $d$ / 95\,\%\,CI not reported in primary publication.%
\footnote{Single-source: P1NP at 4\,h reported as $+153\,\%$ vs.\ baseline for 15\,g and $+53.9\,\%$ for placebo; exact source tables not independently verified.}
&
\textbf{S} (biomarker P1NP; in-vitro ligament proxy; no tendon structural or clinical outcome)
&
Low
&
Baar lab; gelatin donated by industry; co-author holds related patents. Flag.\\

\addlinespace

Lis \& Baar (2019).
\textit{Int J Sport Nutr Exerc Metab} 29(5):526--531. \cite{LB19}
&
Randomised double-blind crossover; n=10 recreationally active males.
&
15\,g vitamin-C-enriched gelatin, 15\,g vitamin-C-enriched hydrolysed collagen, gummy (gelatin+HC), or placebo; 1\,h before 6\,min jump rope.
&
Serum P1NP at 4\,h; serum amino-acid profile.
&
Gelatin and HC both: ${\approx}20\,\%$ rise in P1NP from baseline (trend), not statistically significant vs.\ placebo due to high inter-individual variability. ES and 95\,\%\,CI not reported.
&
\textbf{S} (biomarker P1NP; amino-acid kinetics; no structural or clinical endpoint)
&
Low
&
Same lab as Shaw (2017); no explicit conflict statement.\\

\addlinespace

Praet et al.\ (2019).
\textit{Nutrients} 11(1):76. \cite{PPW+19}
&
Double-blind placebo-controlled crossover pilot RCT; n=20 chronic mid-portion Achilles tendinopathy patients; 6-month protocol.
&
5\,g/day specific collagen peptides (TENDOFORTE\textsuperscript{\textregistered}) $+$ twice-daily calf strengthening vs.\ placebo $+$ strengthening.
&
VISA-A score (pain/function) at 3\,months; tendon microvascularity (Doppler).
&
VISA-A: collagen-first $+12.6$ (95\,\%\,CI 9.7--15.5) vs.\ placebo-first $+5.3$ (2.3--8.3). MCID for VISA-A ${\approx}6.5$. Microvascularity decreased similarly in both groups (no between-group difference).
&
\textbf{F} (VISA-A: validated pain/function scale; clinically meaningful threshold exceeded); \textbf{S} for microvascularity
&
Moderate (pilot)
&
\textbf{Industry-funded (GELITA\,AG / TENDOFORTE\textsuperscript{\textregistered}).} Flag.\\

\addlinespace

Khatri et al.\ (2021).
\textit{Amino Acids} 53:1647--1663. \cite{KNC+21}
&
Systematic review of 15 RCTs; healthy adults and clinical populations.
&
Collagen peptides 5--15\,g/day $+$ exercise, 3--12\,weeks.
&
Collagen synthesis markers, joint pain, body composition.
&
Pooled conclusion: collagen peptides $+$ exercise improve joint pain/function and increase collagen synthesis markers; no consistent additional effect on muscle protein synthesis vs.\ higher-quality proteins. Heterogeneous study designs; pooled ES not calculated.
&
Mixed: \textbf{S} (synthesis markers); \textbf{F} (joint pain/function outcomes in clinical populations)
&
Moderate (systematic review; heterogeneous risk of bias)
&
Some constituent trials industry-funded; authors note this.\\

\addlinespace

Buchalski et al.\ (2026).
\textit{J Funct Morphol Kinesiol} 11(1):130. \cite{BJAL26}%
\footnote{Cited by two independent sources; DOI should be verified via resolver before submission.}
&
Systematic review; 8 RCTs; n=257 (246M:11F); age 18--52\,yr; 3--15\,wk training.
&
Collagen peptides 5--30\,g/day $+$ ${\geq}50$\,mg vitamin~C $+$ high-load training.
&
Tendon cross-sectional area (CSA); stiffness; Young's modulus.
&
3/4 studies: CSA $\uparrow$ significantly vs.\ placebo (collagen: $+11\,\%$ vs.\ $+4.7\,\%$). 2/2 high-dose studies: stiffness and modulus $\uparrow$ significantly. No additive effect on muscular strength. GRADE A for structural outcomes.
&
\textbf{S-structural} (tendon CSA, stiffness, modulus via imaging/mechanics; not injury incidence or climbing performance; large tendons only)
&
Moderate (GRADE A for structural)
&
Not declared.\\

\addlinespace

Nulty et al.\ (2024).
\textit{Am J Physiol Endocrinol Metab}. \cite{NTD+24}%
\footnote{Single-source claim. Requires independent verification.}
&
Randomised crossover; n=8 resistance-trained men, age $49{\pm}8$\,yr.
&
0, 15, or 30\,g vitamin-C-enriched hydrolysed collagen 1\,h pre-resistance exercise.
&
P1NP AUC 0--6\,h.
&
0\,g: $96{\pm}23$\,\textmu g/mL${\times}$h; 15\,g: $134{\pm}23$ ($+40\,\%$, $p{=}0.024$); 30\,g: $169{\pm}28$ ($+76\,\%$ vs.\ 0\,g, $p{<}0.001$; $+26\,\%$ vs.\ 15\,g, $p{=}0.033$).
&
\textbf{S} (biomarker P1NP only; no structural or functional endpoint)
&
Low (small n; middle-aged; surrogate)
&
Not declared.\\

\addlinespace

Sas-Nowosielski et al.\ (2021).
\textit{Int J Environ Res Public Health} 18(10):5370. \cite{SNWK21}
&
Double-blind placebo-controlled; n=15 elite climbers (13M:2F, 1 excluded after injury); 4-week intervention.
&
$\beta$-alanine 4.0\,g/day vs.\ placebo; divided doses.
&
Campus-board intermittent reaches; easy traverse moves-to-failure and time-to-failure.
&
Campus reaches: $d{=}1.55$, $p{=}0.002$; easy traverse moves $+51\,\%$ ($d{=}1.73$); time-to-failure $+32$\,s ($d{=}1.44$, $p{=}0.044$). Hard traverse: no significant benefit.
&
Mixed: campus and traverse tasks are \textbf{S-functional} (climbing-specific but not grade or route completion); no on-route performance endpoint
&
Moderate (climbing-specific; small n; multiple endpoints)
&
None declared.\\

\addlinespace

Sandford et al.\ (2018).
\textit{BMJ Open Sport Exerc Med} 4(1):e000414. \cite{SSWL18}
&
Multicentre double-blind RCT; n=73 rotator-cuff-related shoulder pain; 8\,weeks + exercise.
&
9 capsules/day MaxEPA (1.53\,g EPA $+$ 1.04\,g DHA) vs.\ oleic acid placebo.
&
Oxford Shoulder Score (OSS) at 2 months; SPADI at 3 months.
&
OSS: mean difference $-0.1$ (95\,\%\,CI $-2.6$ to $2.5$; $p{=}0.95$) --- no effect. SPADI: $-8.3$ (95\,\%\,CI $-15.6$ to $-0.94$; $p{=}0.03$) at 3 months favoring omega-3.
&
\textbf{F} (validated pain/disability scores; shoulder not finger; null primary result)
&
Good (adequate n; multi-centre)
&
None declared.\\

\addlinespace

Grgic et al.\ (2022).
\textit{Clin Nutr ESPEN} 47:89--95. \cite{Grg22}
&
Systematic review and meta-analysis; 16 studies; n=353 (34 female).
&
Caffeine 1--7\,mg/kg vs.\ placebo; acute doses.
&
Isometric handgrip strength.
&
$d{=}0.17$ (95\,\%\,CI 0.10--0.23; $p{<}0.001$) overall. Subgroup 1--3\,mg/kg: $d{=}0.20$ (95\,\%\,CI 0.10--0.30; $p{<}0.001$). Small but significant ergogenic effect.
&
\textbf{S} (dynamometric grip strength; no climbing performance outcome; general athletic population)
&
Moderate (some risk of bias; male-predominant)
&
None declared.\\

\addlinespace

Ruiz-L\'opez et al.\ (2026).
\textit{Nutrients} 18(2):284. \cite{RLMAMR+26}
&
Triple-blind randomised crossover RCT; n=13 male climbers.
&
3\,mg/kg caffeine vs.\ placebo; administered 60\,min pre-testing.
&
Pull-up 1RM and power; pull-up endurance; dead-hang time; maximal dead-hang strength; RFD.
&
All outcomes non-significant: pull-up reps $+3.3\,\%$ (95\,\%\,CI $-3.30$ to $4.37$; $g{=}0.061$); maximal grip $+2.4\,\%$ ($p{=}0.060$; $g{=}0.120$); dead-hang $-7.4\,\%$ ($p{=}0.162$). Effect sizes trivial--small.
&
\textbf{S} (pull-up and dead-hang metrics are climbing-adjacent dynamometric surrogates; no on-wall or grade outcome)
&
Moderate (well-controlled; first climbing-specific caffeine RCT; underpowered)
&
None declared.\\

\end{longtable}

\noindent\textit{Note: No study in Table~\ref{tab:suppstudy} used a climbing-specific primary functional outcome (redpoint grade, on-wall endurance, pulley injury incidence). All climbing-specific evidence tiers are extrapolations.}

%---------------------------------------------------------------------
\subsection{Hydrolyzed Collagen Peptides}
\label{subsec:collagen}

\subsubsection{Mechanism}

Type~I collagen constitutes 65--80\,\% of tendon and ligament dry weight, and is the principal structural protein of finger flexor tendons, annular pulleys (A2, A4), and skin dermis --- all high-stress tissues in rock climbing. Collagen triple-helix formation requires hydroxylation of proline and lysine residues by prolyl-4-hydroxylase and lysyl hydroxylase, both of which are ascorbic-acid-dependent dioxygenases; without adequate vitamin~C, hydroxylation is rate-limited and secreted collagen is thermally unstable and subject to intracellular degradation. \textbf{[Evidence-backed: established biochemistry; Myllyharju, 2003, \textit{Matrix Biol}]}

Hydrolyzed collagen peptides (enzymatically processed to ${\approx}2$--5\,kDa; lower molecular weight than thermal gelatin) are absorbed intact from the gut as Pro-Hyp and Hyp-Gly dipeptides, which accumulate in connective tissue and may act as substrates for \textit{de novo} collagen synthesis and/or as signaling molecules upregulating tenocyte collagen gene expression. \textbf{[Bioplausible: mechanism established in vitro (Shigemura et al., 2009, \textit{J Agric Food Chem}); direct in-vivo mechanistic confirmation in human tendon tissue is lacking]}

Intermittent mechanical loading of tendons during exercise increases local blood flow and tenocyte metabolic activity, creating a window during which circulating amino-acid precursors are preferentially incorporated. The Shaw (2017) protocol exploited this by timing collagen ingestion 1\,h pre-exercise to synchronise peak plasma amino-acid elevation with exercise-induced hyperaemia. \textbf{[Bioplausible: inferred from P1NP kinetics; direct tissue-incorporation measurements absent]}

\subsubsection{Shaw et al.\ (2017) --- Critical Appraisal}

Shaw et al.\ (\textit{Am J Clin Nutr} 105:136--143; \cite{SLBR+17}) used a randomised, double-blind, crossover design in n=8 healthy males (age $21{\pm}0.5$\,yr; no tendon pathology). Three conditions --- 5\,g gelatin\,$+$\,48\,mg vitamin~C, 15\,g gelatin\,$+$\,48\,mg vitamin~C, and placebo --- were each administered as a single dose 1\,h before 6\,min rope-skipping, repeated 3$\times$/day for 3 days, with a 7-day washout.

\textbf{Primary outcomes}: (1) Serum P1NP --- a systemic surrogate biomarker of type~I collagen synthesis rate, \textit{not} a direct measure of tendon or pulley collagen content; (2) mechanical properties (stiffness, collagen content) of in-vitro engineered ligament constructs incubated with participant serum --- a cell-culture functional proxy, not a measurement in the participants' own tendons.

\textbf{Results}: The 15\,g condition produced ${\approx}2{\times}$ greater P1NP AUC vs.\ placebo over 4\,h post-exercise ($p{<}0.05$). Engineered ligament stiffness and collagen content were significantly greater with post-15\,g serum vs.\ post-placebo serum. The 5\,g condition was intermediate. Formal Cohen's $d$ and 95\,\%\,CI were not reported; the 2$\times$ P1NP estimate is a within-person comparison from n=8 with low precision.

\textbf{Critical limitations}:
\begin{enumerate}[noitemsep]
  \item n=8, male-only: underpowered; female hormonal effects on collagen turnover (estrogen upregulates collagen synthesis) unaddressed.
  \item Surrogate endpoint: P1NP reflects systemic whole-body collagen synthesis and does not confirm preferential incorporation into finger pulleys or flexor tendons. No tendon biopsy data exist.
  \item Healthy tendons only: tendinopathic tissue has an altered tenocyte phenotype and remodeling dynamics; results cannot be directly extrapolated to injured A2/A4 pulleys.
  \item Acute single-dose design: no chronic dosing, MRI, ultrasound, or functional load-to-failure outcome was assessed.
  \item Gelatin $\ne$ hydrolyzed collagen peptides: gelatin (thermally denatured) has higher molecular weight fragments than enzymatically hydrolyzed collagen (${\approx}2$--5\,kDa). Lis \& Baar (2019) found comparable amino-acid responses when both were vitamin-C-enriched, suggesting functional equivalence under those conditions.
  \item No climbing-specific outcome: rope-skipping loads Achilles/patellar tendons, not finger flexor pulleys.
  \item Industry conflict: the Baar lab has received collagen-industry support; a co-author holds patents on engineered ligament technology. Not disqualifying, but warrants bias assessment.
\end{enumerate}

\begin{evidencebox}{Collagen Pre-Exercise: P1NP Surrogate Benefit (Shaw 2017) · \evidencebacked · ESS 4}
\textbf{Evidence-backed (surrogate biomarker only):} 15\,g gelatin $+$ 48\,mg vitamin~C taken 1\,h before intermittent loading acutely elevates systemic P1NP in healthy young males (${\approx}2{\times}$; $p{<}0.05$) and improves in-vitro engineered ligament mechanics. \textit{Endpoint: surrogate (P1NP biomarker; ex-vivo ligament stiffness). No direct human tendon structural outcome, no injury incidence, no return-to-climbing, no climbing performance endpoint. This does not constitute evidence that climbers heal pulleys faster, increase tendon stiffness, or reduce injury incidence.}
\end{evidencebox}

\subsubsection{Lis \& Baar (2019) Confirmation}

Lis \& Baar (\textit{IJSNEM} 29:526--531; \cite{LB19}): n=10 males, crossover; all active collagen conditions (vitamin-C-enriched gelatin, hydrolyzed collagen, gummy) increased circulating amino acids (glycine, proline, hydroxyproline) similarly; P1NP showed a ${\approx}20\,\%$ trend from baseline that did not reach statistical significance due to high variability. The study confirms the amino-acid-kinetic window (peak ${\approx}1$\,h post-ingestion) and suggests functional equivalence between gelatin and hydrolyzed collagen forms when vitamin~C is co-supplemented. It does not establish a statistically reliable P1NP effect at 15\,g hydrolyzed collagen alone; the non-significance limits claims of confirmed replication.

\subsubsection{Praet et al.\ (2019): Clinical Tendinopathy Data}

Praet et al.\ (\textit{Nutrients} 11:76; \cite{PPW+19}) tested 5\,g/day specific collagen peptides (TENDOFORTE\textsuperscript{\textregistered}) $+$ calf-strengthening vs.\ exercise-only in n=20 chronic Achilles tendinopathy patients (pilot, double-blind, crossover RCT, 6 months). VISA-A improved by $+12.6$ (95\,\%\,CI 9.7--15.5) in the collagen-first group vs.\ $+5.3$ (95\,\%\,CI 2.3--8.3) in the placebo-first group; the between-group difference exceeds the MCID of ${\approx}6.5$ points. Tendon microvascularity decreased similarly in both groups.

\textbf{Limitations}: Pilot scale (n=20); crossover design with potential carry-over effects; single commercial peptide brand (TENDOFORTE\textsuperscript{\textregistered}); \textbf{funded by GELITA AG} (industry-funded: flag); Achilles tendon loading differs substantially from finger flexor pulley loading in climbing; extrapolation to A2/A4 pulleys requires a tissue-type and loading-pattern leap.

\begin{bioplausbox}{Collagen Peptides + Loading for Connective Tissue Recovery in Climbers · \bioplausible · ESS 2}
\textbf{Bioplausible} for collagen peptides $+$ structured loading supporting connective tissue recovery in climbers, drawing on Praet et al.\ (2019). The closest available analog involves a different tendon, a lower dose, a different loading protocol, and an industry-funded trial. No direct evidence exists for climbing-specific pulley healing.
\end{bioplausbox}

\subsubsection{Dose}

Shaw (2017) demonstrated a dose-response between 5\,g and 15\,g for P1NP elevation; 15\,g produced the strongest signal. Lis \& Baar (2019) used 15\,g for all active conditions. Khatri et al.\ (2021; systematic review, 15 RCTs) identified 5--15\,g/day as the studied range, with 15\,g consistently associated with the highest collagen-synthesis marker responses. Praet et al.\ (2019) showed clinical efficacy at 5\,g/day in a chronic tendinopathy context with structured loading, suggesting lower doses may be sufficient for chronic daily protocols. No RCT has tested doses $>$15\,g in young healthy athletes; evidence from one single-source study (Nulty et al., 2024)\footnote{Single-source claim --- requires independent verification.} in middle-aged men suggests 30\,g produces $+26\,\%$ greater P1NP AUC than 15\,g ($p{=}0.033$), but this has not been replicated.

\textbf{[Evidence-backed:]} 15\,g is the best-supported acute pre-exercise dose for collagen synthesis biomarker elevation. No evidence that doses $>$15\,g provide additional clinical benefit in young athletes.

\subsubsection{Timing Window}

Shaw (2017) and Lis \& Baar (2019) both timed collagen ingestion 1\,h before exercise to align the plasma amino-acid peak (${\approx}60$\,min post-ingestion for 15\,g) with exercise-induced tendon hyperaemia. Gastric emptying variability (affected by co-ingested food, hydration, and sympathetic tone) could shift the peak by ${\pm}15$--30\,min. Taking collagen 45--75\,min pre-session adequately approximates the Shaw timing window. Immediate pre-exercise ingestion ($<$15\,min) likely misaligns the amino-acid peak with early high-strain loading.

\textbf{[Bioplausible:]} 1\,h pre-exercise is mechanistically optimal; direct evidence that the 1\,h window is superior to, e.g., 45\,min or 90\,min, does not exist from RCT comparisons.

\subsubsection{The Vitamin~C Co-Ingestion Gap}
\label{subsubsec:vitc}

\begin{warningbox}{\textbf{Protocol Gap: Vitamin\textasciitilde{}C Absent from Collagen Shake}}
The user takes 15\,g hydrolyzed collagen peptides pre-session \textbf{without} vitamin~C co-ingestion. Shaw et al.\ (2017) and Lis \& Baar (2019) both used vitamin-C-enriched preparations (48\,mg ascorbic acid per dose). Without this co-factor, the evidence-backed biomarker benefit \textbf{does not cleanly transfer}.
\end{warningbox}

\textbf{Biochemistry}: Prolyl-4-hydroxylase requires ascorbic acid as an obligatory electron donor; in its absence the hydroxylation cycle stalls, and hydroxylation efficiency falls in proportion to ascorbate depletion. Under frank deficiency ($<$10\,mg/day), collagen triple-helix formation fails entirely (scurvy). At suboptimal non-scorbutic intakes (plasma ascorbate $<$23\,$\mu$mol/L), hydroxylation efficiency is proportionally reduced. \textbf{[Evidence-backed: established enzymology]}

\textbf{Shaw protocol context}: 48\,mg vitamin~C was co-ingested per dose. The US RDA for vitamin~C is 90\,mg/day for men and 75\,mg/day for women; 48\,mg is sub-RDA but was not intended as a pharmacological dose --- its purpose was to standardise local ascorbate availability during the post-ingestion window without confounding the collagen signal.

\textbf{Scenario analysis} for the user's protocol:

\begin{center}
\begin{tabular}{>{\raggedright\arraybackslash}p{3.2cm}
                >{\raggedright\arraybackslash}p{3.8cm}
                >{\raggedright\arraybackslash}p{5.5cm}}
\toprule
\textbf{Vitamin~C status} & \textbf{Plasma ascorbate} & \textbf{Expected collagen synthesis response} \\
\midrule
Adequate habitual diet (${\geq}5$ fruit/veg servings/day; ${\geq}100$\,mg/day dietary C) & ${\geq}50$\,$\mu$mol/L (replete) & Prolyl hydroxylase likely not substrate-limited; response may approximate Shaw protocol. Co-ingestion provides marginal additional ascorbate. \\
\addlinespace
Suboptimal intake (common in athletes in caloric restriction; $<$50\,mg/day dietary C) & ${<}23$\,$\mu$mol/L (sub-optimal) & Substrate (Pro, Gly, Hyp) elevated but hydroxylation rate-limited; collagen synthesis efficiency impaired; evidence-backed benefit does not apply. \\
\bottomrule
\end{tabular}
\end{center}

\textbf{Critical gap}: No RCT has directly compared 15\,g collagen \textit{with} vs.\ \textit{without} co-ingested vitamin~C under controlled dietary ascorbate conditions. The claim that omitting vitamin~C definitively abolishes benefit is \textbf{[Bioplausible]}, not experimentally proven. However, the claim that the user's protocol replicates the Shaw design is \textbf{false}: the designs differ in a mechanistically important cofactor. The uncertainty is asymmetric --- co-ingesting 50--100\,mg vitamin~C has no downside (cost $<$\$0.01/day; no risk at this dose) and eliminates the bottleneck. \textbf{Absence of co-ingestion is a correctable gap with no counterargument.}

\textbf{[Bioplausible]:} Adding ${\approx}50$\,mg ascorbic acid to the collagen shake is the single highest-priority protocol modification for the user. Whether it is strictly necessary depends on habitual dietary vitamin~C status, which has not been measured.

\subsubsection{Claim: Collagen for Climbing Skin Recovery}

\textbf{Claim}: Oral hydrolyzed collagen peptides accelerate healing of flappers (partial-thickness skin avulsions), skin tears, and chronic fingertip skin thinning in climbers.

\textbf{Evidence}: Puškarić et al.\ (2024; \textit{Nutrients}, \cite{FLAGA+24})\footnote{\toverify: only two sources cite a skin-specific RCT; the DOI cited by other sources resolves to Fernández-Lázaro 2024 omega-3 review. The skin RCT is a single-source claim requiring independent DOI verification.} examined collagen (5\,g $+$ 80\,mg vitamin~C/day $\times$ 16\,wk) in aging women (n=87), showing improved dermis density and texture vs.\ placebo. This population (cosmetic aging) and outcome (dermal density) differ mechanistically from acute traumatic skin avulsion and callus reformation in climbers.

No peer-reviewed RCT has tested oral collagen peptides for climbing-specific skin injuries.

\begin{folklorebox}{Collagen Accelerates Healing of Climbing Flappers and Split Tips · \folkloreverified · ESS 1}
\textbf{Folklore:} ``Collagen accelerates healing of climbing flappers and split tips.'' Source: climbing community forums (r/climbharder), coach blogs (Hooper's Beta, Lattice Training social media), supplement marketing. No peer-reviewed climbing-specific or acute wound-healing evidence exists for this claim. Bioplausible substrate mechanism (collagen provision $\to$ dermal fibroblast activity) does not constitute evidence.
\end{folklorebox}

\subsubsection{Claim: Collagen for Pulley and Tendon Remodeling Under Active Climbing Load}

\textbf{Mechanism} \textbf{[Bioplausible]}: Tendons and annular pulleys are ${\approx}70\,\%$ type~I collagen by dry weight. During loading, tenocytes upregulate TGF-$\beta$1 and IGF-1 signaling, increasing collagen turnover. Net anabolic balance requires adequate substrate (proline, glycine, hydroxyproline) and co-factors (vitamin~C). Provision of 15\,g hydrolyzed collagen elevates circulating Gly ($+376$\,$\mu$mol/L), Pro ($+162$\,$\mu$mol/L), and Hyp ($+105$\,$\mu$mol/L) at ${\approx}1$\,h post-ingestion (Shaw, 2017); active climbing concurrently increases tendon perfusion. The anabolic-window model is mechanistically coherent. \textbf{[Bioplausible: mechanism plausible; no climbing RCT measures pulley or flexor tendon collagen content change attributable to oral supplementation]}

Buchalski et al.\ (2026; systematic review) report that collagen 15--30\,g/day $+$ vitamin~C $+$ high-load training (${\geq}70\,\%$ 1RM resistance) increases tendon CSA ($+11\,\%$ vs.\ $+4.7\,\%$ placebo) and stiffness/modulus in 3/4 studies; GRADE A for structural outcomes.\footnote{Two independent sources cite Buchalski et al.\ 2026. Source metadata requires resolver verification before submission.} This is the strongest structural-outcome evidence available, but it involves large tendons under resistance-training load, not small-diameter finger pulleys under climbing-specific isometric crimp load.

\begin{bioplausbox}{Collagen for Pulley and Tendon Remodeling Under Climbing Load · \bioplausible · ESS 2}
\textbf{Bioplausible:} Collagen peptides $+$ vitamin~C $+$ structured loading may support type~I collagen anabolism in finger flexor tendons and A2/A4 pulleys. The Praet (2019) Achilles tendinopathy data and Buchalski (2026) structural data are the closest available analogs. No climbing-specific RCT exists for pulley healing, injury prevention, or return-to-climbing outcomes.
\end{bioplausbox}

\subsubsection{User Protocol Assessment and Recommended Correction}

\textbf{Current protocol}: 15\,g hydrolyzed collagen peptides pre-climbing session, no vitamin~C, unspecified exact timing.

\begin{center}
\begin{tabular}{>{\raggedright\arraybackslash}p{3.5cm}
                >{\raggedright\arraybackslash}p{2.5cm}
                >{\raggedright\arraybackslash}p{3.2cm}
                >{\raggedright\arraybackslash}p{3.5cm}}
\toprule
\textbf{Parameter} & \textbf{User protocol} & \textbf{Evidence-backed protocol} & \textbf{Gap} \\
\midrule
Dose & 15\,g \checkmark & 15\,g (Shaw/Lis/Khatri) & None \\
\addlinespace
Vitamin~C & \textbf{None} & ${\geq}48$--50\,mg co-ingested & \textbf{Critical gap} \\
\addlinespace
Timing & ``Before climbing'' (unspecified) & 45--75\,min pre-session & Clarify; risk if $<$15\,min \\
\addlinespace
Frequency & Session days only & Daily (Shaw: 3$\times$/day $\times$3\,d; Baar 2017 recommendation: daily) & Uncertain gap; rest-day dosing untested \\
\addlinespace
Loading stimulus & Rock climbing & Intermittent loading ${\geq}6$\,min at moderate--high intensity & Climbing provides adequate stimulus \checkmark \\
\addlinespace
Monitoring & Unknown & Periodic grip/pain VAS assessment & No monitoring \\
\bottomrule
\end{tabular}
\end{center}

\textbf{Minimum corrective action}: Co-ingest 50--100\,mg ascorbic acid (e.g., 100\,mg tablet or a small glass of citrus juice) with the 15\,g collagen shake, taken 45--60\,min before session start. This aligns the protocol with Shaw et al.\ (2017) and eliminates the mechanistic bottleneck at prolyl hydroxylase.

%---------------------------------------------------------------------
\subsection{Creatine Monohydrate}
\label{subsec:creatine}

\textbf{Mechanism}: Creatine supplementation saturates intramuscular phosphocreatine (PCr) stores, accelerating ATP resynthesis during high-intensity efforts of $\leq$15\,s (PCr-dependent energy system). Secondary: creatine may augment muscle hypertrophy via cell volumisation and mTOR pathway activation. Relevant to climbing: maximal boulder moves, contact-strength efforts, campus board power, and repeated attempts with brief rest intervals. \textbf{[Evidence-backed: well-replicated mechanism]}

\textbf{Evidence}: Multiple meta-analyses document creatine's efficacy for muscular strength and power in resistance-trained individuals. No climbing-specific RCT has measured creatine effects on boulder problem performance, maximum hang time, or campus board output in trained climbers. \textbf{[Evidence-backed for general strength/power; Bioplausible for climbing]}

\textbf{Dose}: No loading phase required for saturation if 3--5\,g/day maintained for $\geq$28 days. If rapid saturation is desired: 20\,g/day $\div$ 4 doses $\times$ 5--7 days, then 3--5\,g/day maintenance. Creatine monohydrate is the only form with robust evidence; effervescent, ethyl ester, and buffered variants show no superiority.

\textbf{Climbing-specific caveat}: Body mass increases of ${\approx}$0.5--2\,kg from intracellular water retention may impair strength-to-weight ratio and thus on-wall performance, particularly in weight-sensitive disciplines (lead climbing, slabs, long traverses). The net effect for power-oriented bouldering disciplines is likely positive, but this has not been tested in a controlled climbing trial.

\begin{evidencebox}{Creatine: Strength and Power in Resistance-Trained Athletes · \evidencebacked · ESS 4}
\textbf{Evidence-backed} for general strength and power in resistance-trained athletes (multiple meta-analyses). \textbf{Bioplausible} for climbing performance improvement; body-mass caveat is real and applies specifically to climbing.
\end{evidencebox}

\begin{toverifybox}{Creatine Body-Mass Penalty in Climbing Performance · \toverify · ESS 1}
\textbf{\toverify:} Community claims that ``creatine makes you too heavy to climb'' are anecdotal and not supported by any controlled climbing-performance study. The magnitude of mass gain (0.5--2\,kg) and its effect on climbing-specific metrics remain empirically untested. Only two sources identified (r/climbharder; commercial supplement blogs); requires $\geq$3 independent climbing-community sources for \folkloreverified status.
\end{toverifybox}

%---------------------------------------------------------------------
\subsection{Omega-3 Fatty Acids (EPA/DHA)}
\label{subsec:omega3}

\textbf{Mechanism}: EPA and DHA compete with arachidonic acid for COX-2 and LOX enzymatic pathways, shifting eicosanoid and leukotriene profiles toward less pro-inflammatory mediators. Potential relevance to climbing: modulating tendon/ligament inflammation in overuse pathology (pulley tendinopathy, medial epicondylopathy), reducing training-induced DOMS, and supporting synovial joint health in finger DIP/PIP joints. \textbf{[Evidence-backed: biochemical mechanism; Calder, 2017, \textit{Eur J Clin Nutr}]}

\textbf{Best tendon-specific RCT evidence}: Sandford et al.\ (2018; \cite{SSWL18}), n=73 rotator-cuff-related shoulder pain, 2.57\,g/day EPA+DHA $\times$ 8 weeks $+$ exercise: no clinically meaningful OSS benefit at 2 months (mean difference $-0.1$, 95\,\%\,CI $-2.6$ to $2.5$; $p{=}0.95$); modest SPADI improvement at 3 months ($-8.3$, 95\,\%\,CI $-15.6$ to $-0.94$; $p{=}0.03$) of uncertain clinical significance. No direct benefit for structural tendon outcomes. \textbf{[Evidence-backed for systemic anti-inflammatory effect; Evidence-backed (null result) for tendinopathy structural/pain outcomes in this population]}

A Mendelian randomisation study found genetically higher omega-3 levels were associated with an $+18\,\%$ increase in Achilles tendinitis risk.\footnote{Single-source claim --- requires independent verification. Mendelian randomisation cannot establish causality and is subject to pleiotropy; interpret with caution.}

\textbf{Dose}: 2--3\,g/day combined EPA+DHA is the dose range used in most anti-inflammatory trials. No dose-response curve for tendon outcomes. Daily supplementation; timing is not acutely critical. Antiplatelet effect at $>$3\,g/day; anticoagulant interaction at high doses; fish oil peroxidation if stored improperly; algal sources available for vegetarians.

\begin{bioplausbox}{Omega-3 for Tendon Inflammation Modulation in Climbers · \bioplausible · ESS 3}
\textbf{Bioplausible} for modulation of tendon-related inflammation in climbers with overuse pathology. Current controlled evidence does not confirm clinically meaningful tendon-structural or injury-prevention benefit. The best available tendinopathy RCT (Sandford 2018, rotator cuff) shows no primary-outcome benefit. No climbing-specific RCT exists.
\end{bioplausbox}

\begin{toverifybox}{Omega-3 Strengthens Tendons or Prevents Pulley Injuries in Climbers · \toverify · ESS 0}
\textbf{\toverify:} Claims that omega-3 ``strengthens tendons'' or ``prevents pulley injuries'' in climbers. Only two sources identified (physiotherapy blogs; sports nutrition marketing); requires $\geq$3 independent climbing-community sources. Not supported by current controlled evidence (Sandford 2018 null result).
\end{toverifybox}

%---------------------------------------------------------------------
\subsection{Magnesium}
\label{subsec:magnesium}

\textbf{Mechanism}: Mg$^{2+}$ is cofactor for $>$300 enzymatic reactions, including Mg-ATP complex formation, muscle relaxation (antagonises Ca$^{2+}$ at the sarcoplasmic reticulum), and sleep architecture modulation (GABA-receptor agonism). Relevant to climbing: neuromuscular recovery between moves, sleep quality for training adaptation, and potentially cramp prevention. \textbf{[Evidence-backed: physiological role]}

\textbf{Evidence}: Magnesium supplementation corrects documented deficiency and may improve performance in deficient individuals; effects in magnesium-replete athletes are small or absent. Cochrane review and a JAMA RCT (n=94) found magnesium oxide no better than placebo for nocturnal leg cramps ($\Delta -0.38$ cramps/week in both groups).\footnote{Single-source claim for the JAMA cramp RCT --- requires independent verification.} For sleep, a specific magnesium bisglycinate (250\,mg elemental/day) trial reported Insomnia Severity Index improvement ($d{\approx}0.2$, CI reported) at 4 weeks in poor sleepers.\footnote{Single-source claim --- requires independent verification.} No climbing-specific RCT exists for magnesium effects on forearm pump, grip endurance, or training recovery. \textbf{[Bioplausible for climbing performance benefit in athletes with suboptimal status]}

Climbers in caloric restriction, heavy sweating, and high training loads are at elevated risk of suboptimal (not frank deficiency) magnesium status. Serum magnesium is an insensitive biomarker; erythrocyte or ionised magnesium testing is more informative.

\textbf{Dose}: 200--400\,mg elemental magnesium/day (as glycinate, malate, or citrate; superior bioavailability vs.\ oxide). Evening dosing to support sleep benefit. Diarrhoea common at $>$400\,mg oxide.

\begin{bioplausbox}{Magnesium for Sleep and Neuromuscular Recovery in Suboptimal-Status Climbers · \bioplausible · ESS 2}
\textbf{Bioplausible} for sleep-mediated recovery benefit in climbers with suboptimal magnesium status. \textbf{Evidence-backed} for correcting documented deficiency. \textbf{Evidence-backed (null)} for cramp prevention in general populations.
\end{bioplausbox}

\begin{toverifybox}{Magnesium Prevents Forearm Pump or Finger Cramps in Climbers · \toverify · ESS 0}
\textbf{\toverify:} Claims that magnesium ``prevents forearm pump'' or ``stops finger cramps on routes.'' Only two sources identified (coaching forums; gym culture); requires $\geq$3 independent climbing-community sources. No targeted trial supports this in euvolaemic, well-fed climbers.
\end{toverifybox}

%---------------------------------------------------------------------
\subsection{Beta-Alanine}
\label{subsec:betaalanine}

\textbf{Mechanism}: $\beta$-alanine is the rate-limiting precursor for muscle carnosine synthesis. Carnosine ($\beta$-alanyl-L-histidine; pKa ${\approx}$6.83) acts as an intramuscular H$^+$ buffer, attenuating acidosis during high-intensity glycolytic work. Forearm pump in rock climbing --- the accumulating burning sensation from sustained gripping on crux sections and power-endurance routes --- is mechanistically attributable to intramuscular H$^+$ accumulation from glycolysis; carnosine's buffering capacity targets precisely this mechanism. \textbf{[Bioplausible for climbing; mechanism Evidence-backed in general contexts]}

\textbf{Evidence}: Hobson et al.\ (2012, \textit{Amino Acids}, meta-analysis, k=15 RCTs): $\beta$-alanine significantly increased exercise capacity for efforts of 60--240\,s (ES = 0.374, 95\,\%\,CI 0.140--0.607). Trexler et al.\ (2015, \textit{J Int Soc Sports Nutr}, meta-analysis, k=57 studies): mean performance improvement $+2.85\,\%$ (95\,\%\,CI 0.37--10.49); Hedges $g = 0.374$ for 60--240\,s efforts specifically. This duration range (1--4\,min) corresponds to sustained pump sections and power-endurance sequences on sport climbing routes. \textbf{[Evidence-backed for 60--240\,s high-intensity exercise]}

\textbf{Climbing-specific RCT}: Sas-Nowosielski et al.\ (2021; \cite{SNWK21}): $\beta$-alanine 4\,g/day $\times$ 4 weeks in n=15 elite climbers (double-blind, placebo-controlled). Campus-board reaches: $d{=}1.55$, $p{=}0.002$; easy traverse moves: $+51\,\%$, $d{=}1.73$; easy traverse time-to-failure: $+32$\,s, $d{=}1.44$. Hard traverse: no significant benefit. The large effect sizes should be interpreted cautiously given the small sample (n=15, with 1 excluded) and multiple endpoints without correction for multiplicity. \textbf{[Evidence-backed for specific climbing-relevant endurance tasks; small n; moderate quality]}

\textbf{Dose}: 3.2--6.4\,g/day in divided doses (0.8--1.6\,g per dose, every 3--4\,h) to minimise paresthesia. Carnosine loading requires 4--12 weeks at therapeutic doses for ${\approx}$60--80\,\% increase in muscle carnosine content.

\textbf{Caveat}: Paresthesia (harmless tingling, primarily face/scalp/hands) is near-universal at single doses $>$800\,mg; divided or slow-release formulations mitigate this. Performance benefit is most pronounced for efforts of 60--240\,s; submaximal long-duration routes with low anaerobic contribution may see less benefit. No benefit for maximal strength or power output.

\begin{evidencebox}{Beta-Alanine: Climbing Endurance Task Benefit (Sas-Nowosielski 2021) · \evidencebacked · ESS 7}
\textbf{Evidence-backed} for 60--240\,s high-intensity exercise (multiple meta-analyses, Hedges $g = 0.374$) and for specific climbing-endurance tasks in elite climbers (campus board, easy traverse; Sas-Nowosielski 2021, $d = 1.44$--1.73, but n=15 only). Dose: 3.2--6.4\,g/day divided $\times \geq$4 weeks. \textit{Endpoint: campus-board and traverse tasks are climbing-specific surrogate-functional hybrids; no redpoint grade, route completion rate, or injury incidence endpoint. Hard traverse showed no benefit.}
\end{evidencebox}

\begin{toverifybox}{Beta-Alanine Eliminates Forearm Pump in Climbing · \toverify · ESS 0}
\textbf{\toverify:} Claims that $\beta$-alanine ``eliminates pump'' in climbing. Only two sources identified (climbing blogs; nutrition forums); requires $\geq$3 independent climbing-community sources. Current data show improvement in specific tasks, not elimination of pump; hard traverse performance showed no benefit in the only climbing RCT (Sas-Nowosielski 2021).
\end{toverifybox}

%---------------------------------------------------------------------
\subsection{Caffeine}
\label{subsec:caffeine}

\textbf{Mechanism}: Non-selective adenosine receptor antagonist (A1, A2A). Reduces perceived exertion, enhances neuromuscular activation, increases pain tolerance, potentiates motor unit recruitment, and accelerates calcium release from the sarcoplasmic reticulum. Relevant to climbing: sustained isometric grip force, delayed onset of perceived forearm pump, acute power output on difficult moves. \textbf{[Evidence-backed]}

\textbf{Best general evidence}: Grgic et al.\ (2022, \textit{Clin Nutr ESPEN} 47:89--95; \cite{Grg22}): meta-analysis, k=16 studies, n=353; isometric handgrip strength: $d{=}0.17$ (95\,\%\,CI 0.10--0.23; $p{<}0.001$). Subgroup 1--3\,mg/kg: $d{=}0.20$ (95\,\%\,CI 0.10--0.30). Small but statistically significant ergogenic effect on a measurement directly analogous to climbing grip. \textbf{[Evidence-backed for isometric grip strength]}

Grgic et al.\ (2020, \textit{Br J Sports Med} 54:681--688; \cite{WTNZ18}): k=34 studies; muscle strength ES = 0.20 (95\,\%\,CI 0.13--0.26); muscular endurance ES = 0.41 (95\,\%\,CI 0.24--0.58). \textbf{[Evidence-backed for general muscular endurance]}

\textbf{Climbing-specific RCT}: Ruiz-L\'opez et al.\ (2026; \cite{RLMAMR+26}): triple-blind randomised crossover, n=13 male climbers, 3\,mg/kg caffeine vs.\ placebo, 60\,min pre-testing. Pull-up reps: $+3.3\,\%$ (95\,\%\,CI $-3.30$ to $4.37$; $g{=}0.061$); maximal grip: $+2.4\,\%$ ($p{=}0.060$; $g{=}0.120$); dead-hang: $-7.4\,\%$ ($p{=}0.162$). All outcomes non-significant; effect sizes trivial--small. Subjective energy and alertness increased ($p{<}0.05$). This is the only published climbing-specific caffeine RCT and is underpowered (n=13); it does not rule out a small real effect at higher doses. \textbf{[Evidence-backed: 3\,mg/kg does not meaningfully improve proxy climbing performance at this sample size]}

\textbf{Dose}: 3--6\,mg/kg body mass, 45--60\,min pre-session. Lower doses (3\,mg/kg) sufficient for endurance and perceived-exertion benefit with fewer side effects; higher doses (5--6\,mg/kg) may increase maximal strength/power but risk anxiety, tachycardia, and sleep disruption. Habitual consumers show attenuated response; 48--96\,h abstinence before target sessions partially restores the effect.

\textbf{Climbing-specific caveat}: Fear and anxiety are performance factors in climbing, particularly on serious routes and in competition. Caffeine increases sympathetic arousal and may exacerbate anxiety-driven performance decrements, tremor, and hesitation on high-commitment moves. Sleep disruption at doses $>$3\,mg/kg within 6\,h of sleep onset reduces recovery quality.

\begin{evidencebox}{Caffeine: Isometric Grip Strength and Endurance Benefit · \evidencebacked · ESS 4}
\textbf{Evidence-backed} for isometric handgrip strength ($d{=}0.17$, Grgic 2022) and general muscular endurance ($d{=}0.41$, Grgic 2020). \textbf{Evidence-backed (null at 3\,mg/kg, n=13)} for climbing-specific pull-up/grip metrics (Ruiz-L\'opez 2026). Dose 3--6\,mg/kg 45--60\,min pre-session; individual tolerance and habituation are critical modifiers. \textit{Endpoint: dynamometric surrogate (grip strength, endurance metrics); no redpoint grade, route completion, or competition ranking outcome in any trial. The only climbing-specific RCT (Ruiz-L\'opez 2026) showed null results on all proxy climbing metrics.}
\end{evidencebox}

%---------------------------------------------------------------------
\subsection{What the Literature Cannot Currently Resolve}
\label{subsec:suppgaps}

\begin{enumerate}[noitemsep]

  \item \textbf{Collagen $+$ vitamin~C and pulley structure in climbers}: No study has measured A2/A4 pulley or finger flexor tendon CSA, stiffness, or load-to-failure as a function of oral collagen supplementation using high-resolution ultrasound, MRI, or biopsy in any climbing population. All positive structural data (Buchalski 2026) involve large tendons under resistance-training loads; small-diameter, high-strain pulley responses are unknown.

  \item \textbf{Minimum effective vitamin~C co-dose}: Shaw (2017) used 48\,mg; Lis (2019) used enriched preparations. No factorial RCT has compared collagen $\pm$ vitamin~C under controlled dietary ascorbate conditions, so the threshold above which co-ingestion is redundant for athletes with replete diets remains undefined.

  \item \textbf{Dose-response above 15\,g}: No RCT has compared 15\,g vs.\ 20\,g vs.\ 30\,g collagen with functional tendon outcomes in young healthy athletes. The 30\,g data (Nulty 2024, single source, middle-aged men) suggest an additional P1NP increment but have not been independently replicated.

  \item \textbf{Collagen and climbing skin repair}: No controlled trial has assessed time-to-heal for flappers, split tips, or chronic skin thinning in climbers receiving collagen supplementation vs.\ placebo under standardised wound-care conditions.

  \item \textbf{Beta-alanine and perceived pump in climbing}: The Sas-Nowosielski (2021) RCT showed campus board and easy traverse improvements but no hard traverse benefit. Whether carnosine loading reduces perceived pump and improves on-route performance in trained sport climbers (where vascular occlusion, not only metabolic acidosis, contributes to pump) is not established.

  \item \textbf{Caffeine at doses $>$3\,mg/kg in climbing}: The only climbing-specific RCT (Ruiz-L\'opez 2026) used 3\,mg/kg only. Whether 5--6\,mg/kg produces clinically meaningful improvement in actual climbing performance (redpoint grade, route completion percentage, competition ranking) without offsetting anxiety or tremor remains untested in climbing populations.

  \item \textbf{Omega-3 and finger tendon remodeling}: The available tendinopathy RCT (Sandford 2018, rotator cuff) showed no primary-outcome benefit. No RCT has examined EPA/DHA effects specifically on finger flexor tendons or pulleys, or on A2 pulley injury recurrence rate.

  \item \textbf{Interaction effects}: No study has examined whether collagen $+$ vitamin~C interacts with creatine, beta-alanine, or caffeine in terms of tendon adaptation, injury risk, or climbing-specific performance in a factorial or combination design.

\end{enumerate}
